\begin{theorem}[最小状态复杂度可编码停机时间]\label{thm:conclusion-minimal-state-complexity-halting}
存在一可计算映射 \((M,x)\mapsto \pi_{M,x}\)，将任意图灵机 \(M\) 与输入 \(x\) 送到一个全定义的序贯转导 \(\pi_{M,x}\)（在某个正则语言上定义），使得：
\begin{enumerate}
  \item 对每个固定 \((M,x)\)，\(\pi_{M,x}\) 的每个输出纤维大小至多为 \(2\)；
  \item 记 \(T(M,x)\in\NN\cup\{\infty\}\) 为停机时间（不停机记为 \(\infty\)），并令 \(s(\pi_{M,x})\) 为实现 \(\pi_{M,x}\) 的最小确定性有限状态转导器的状态数。
  则当 \(T(M,x)<\infty\) 时有
  \[
  s(\pi_{M,x})=T(M,x)+1,
  \]
  而当 \(T(M,x)=\infty\) 时 \(s(\pi_{M,x})=O(1)\)。
\end{enumerate}
从而映射 \((M,x)\mapsto s(\pi_{M,x})\) 不可计算；并且性质“存在统一常数 \(S\) 使对所有 \((M,x)\) 有 \(s(\pi_{M,x})\le S\)”在该族上不可判定。
\leanverified{paper\_conclusion\_minimal\_state\_complexity\_halting}
\end{theorem}
\begin{proof}
取字母表 \(\Sigma:=\{0,1,\#\}\)，并令正则语言
\(
L:=\{\,1^n\# b:\ n\ge 0,\ b\in\{0,1\}\,\}
\)。
对每个 \((M,x)\)，定义 \(\pi_{M,x}:L\to L\) 为
\[
\pi_{M,x}(1^n\# b)
:=
0^n\#\begin{cases}
b,& M(x)\ \text{在 }n\text{ 步内未停机},\\
0,& M(x)\ \text{在 }n\text{ 步内已停机}.
\end{cases}
\]
谓词“在 \(n\) 步内停机”可由有限步模拟判定，故 \(\pi_{M,x}\) 可计算且全定义。
并且对每个输出 \(0^n\# c\)，其原像至多为 \(1^n\#0\) 与 \(1^n\#1\) 两点，故纤维大小 \(\le 2\)。
\par
若 \(T:=T(M,x)<\infty\)，则当且仅当 \(n\ge T\) 时输出末位被强制为 \(0\)。
对任意 \(0\le i<T\)，考虑前缀 \(1^i\) 与 \(1^{i+1}\)。
它们在共同后缀 \(1^{T-i-1}\#1\) 的续接下，分别对应输入长度 \(T-1\) 与 \(T\)，从而输出末位分别为 \(1\) 与 \(0\)。
因此从 Myhill--Nerode 最小化原理（或等价的“未来输出函数不同 \(\Rightarrow\) 状态必须区分”）可知，
在读入前缀 \(1^i\) 后必须处于彼此不同的状态，\(i=0,1,\dots,T\) 共需至少 \(T+1\) 个状态，故 \(s(\pi_{M,x})\ge T+1\)。
另一方面，用状态计数器 \(\{0,1,\dots,T\}\)（到达 \(T\) 后吸收）即可实现上述阈值行为，故 \(s(\pi_{M,x})\le T+1\)，从而 \(s(\pi_{M,x})=T+1\)。
\par
若 \(T(M,x)=\infty\)，则 \(\pi_{M,x}\) 退化为逐符号转导 \(1\mapsto 0\)、\(\#\mapsto \#\)、\(b\mapsto b\)，显然可由常数状态实现，故 \(s(\pi_{M,x})=O(1)\)。
\par
最后，若 \((M,x)\mapsto s(\pi_{M,x})\) 可计算，则可由 \(s(\pi_{M,x})\) 判定 \(T(M,x)\) 是否有限并恢复其值，与停机问题不可判定性矛盾。证毕。
\end{proof}

\begin{corollary}[最小寄存器数的不可计算性：弱硬件判据已足以编码停机性]\label{cor:conclusion-min-register-halting}
固定一种标准同步序贯电路模型：以 NAND 作为唯一组合逻辑门，并以 \(k\) 个二值寄存器（触发器）作为状态记忆。
对任意序贯转导 \(\pi\) 定义
\[
k_{\min}(\pi):=\min\Bigl\{k:\ \exists\ \text{上述模型的电路实现 }\pi\Bigr\}.
\]
则映射 \((M,x)\mapsto k_{\min}(\pi_{M,x})\) 不可计算。
并且存在一个全定义可计算的前置变换 \((M,x)\mapsto(M^{(+1)},x)\)（先执行一步空操作再模拟 \(M\)），使得
\[
k_{\min}(\pi_{M^{(+1)},x})=
\begin{cases}
0,&T(M,x)=\infty,\\
\ge 1,&T(M,x)<\infty,
\end{cases}
\]
从而“是否需要至少一个记忆触发器”这一极弱判据已足以编码停机性。
\leanverified{paper\_conclusion\_min\_register\_halting}
\end{corollary}

\begin{proof}
任何含 \(k\) 个二值寄存器的确定性同步序贯电路，其可达内部状态数至多为 \(2^k\)。
因此对任意 \(\pi\) 有下界
\(
k_{\min}(\pi)\ge \lceil\log_2 s(\pi)\rceil
\)。
由定理 \ref{thm:conclusion-minimal-state-complexity-halting}，当 \(T(M,x)=T<\infty\) 时 \(s(\pi_{M,x})=T+1\)，从而 \(k_{\min}(\pi_{M,x})\ge \lceil\log_2(T+1)\rceil\)；
而当 \(T(M,x)=\infty\) 时 \(\pi_{M,x}\) 退化为常数状态可实现的逐符号转导，故 \(k_{\min}(\pi_{M,x})=0\)。
对加一填充 \(M^{(+1)}\)，若 \(T(M,x)<\infty\) 则 \(T(M^{(+1)},x)\ge 1\)，从而 \(s(\pi_{M^{(+1)},x})\ge 2\) 并强制 \(k_{\min}\ge 1\)；
若 \(T(M,x)=\infty\) 则仍保持 \(k_{\min}=0\)。
因此若 \((M,x)\mapsto k_{\min}(\pi_{M,x})\) 可计算，即可判定停机性，矛盾。证毕。
\end{proof}

\endinput
